\begin{yl}{5}{Kolmanda mõõtme saladus}{portaal}{1 sek / 3 sek}{100 punkti}
  \emph{Idee, teostus ja lahenduse selgitus: Tähvend Uustalu}

  Mis eristab seda ülesannet tüüpilisest ``labürindiülesandest'', mida saab
  ``tavalise'' BFS või DFS abil lahendada? Kuna portaalid
  mõjutavad ka münte, siis võib münt ``eest ära liikuda''. Vaatleme näiteks alloleval
  joonisel kujutatud olukorda. Kui läheme portaali 1 juurde ja tõmbame kangist,
  siis liigume küll portaali 2, kuid münt liigub samal ajal portaali 3. Kui
  uuesti kangist tõmmata, jõuame portaali 3 juurde, kus münt enne oli, aga münt
  liigub samal ajal portaali 4 ja nii edasi. Münti ei olegi võimalik kätte saada.

  \begin{center}
    \includegraphics[scale=1.5]{example1}
  \end{center}

  Kui aga osade portaalide vahel oleks võimalik liikuda, siis saab kavalalt liikudes
  mündi enda kätte manipuleerida. Paremal pool olevas näites saab mündi kätte nii:
  kõndida portaali 1 juurde, tõmmata kangist, kõndida portaali 4 juurde, tõmmata 3
  korda kangist. Seejärel uuesti kõndida portaali 4 juurde ja tõmmata 3 korda kangist.
  Lõpuks oleme portaali 2 juures. Münt on meie seitsme kangitõmbega liikunud
  $2 \to 3 \to 4 \to 5 \to 1 \to 2 \to 3 \to 4$. Saame kõndida portaali 4 juurde
  ja mündi alati üles võtta.

  Ent mitte alati ei ole portaalide vahel kõndimisest abi. Allolevas näites
  ei ole jälle münti võimalik üles korjata.

  \begin{center}
    \includegraphics[scale=1.5]{example2}
  \end{center}

  Seega ei olegi nii lihtne välja selgitada, milliseid münte on üles võimalik
  korjata. Väikeste $K$ korral (näiteks $K \le 4$) võib iseenesest kõik
  variandid kasvõi paberil läbi mõelda. Suuremates testides on vaja süsteemsemat
  lähenemist.

  Esiteks paneme tähele, et kui piisavalt mitu korda järjest kangist tõmmata,
  jõuab kõik lõpuks algseisu tagasi. Kindlasti juhtub see $K! = 1 \cdot 2 \cdot 3 \cdots K$
  kangitõmbe järel, kuigi see võib juhtuda ka varem.\footnote{See juhtub kindlasti
    alati ka $\text{VÜK}(1, 2, \ldots, K)$ kangitõmbe järel, aga sõltuvalt
    konkreetsest permutatsioonist ka veelgi varem.} Seega, kui tõmmata kangist
  $K! - 1$ korda, liiguvad iga $i$ kohta kõik portaali $A[i]$ juures olevad objektid
  portaali $i$. Nii on meil võimalik ``$-1$ korda'' kangist tõmmata.

  See tähendab, et kõiki käike on võimalik ``tagasi võtta'': kui astume naaberruudule, võime
  tagasi eelmisele ruudule astuda; kui tõmbame kangist, võime $-1$ korda kangist tõmmata.
  Seega võime mõelda igast mündist individuaalselt: kui valime mingi mündi ja selle
  mündini on üleüldse võimalik liikuda, siis liigume mingil viisil selle mündini,
  korjame üles ja läheme mööda tuldud teed alguspunkti tagasi. Kõik ülejäänud
  mündid on siis täpselt samadel positsioonidel, nagu alguses. Seega piisab meil
  ülesande lahendamiseks sellest, kui arvutame iga mündi kohta eraldi välja, kas seda on võimalik
  üles korjata: müntide kogumise järjekord ei loe ja kunagi ei teki olukorda, kus peame
  kahe mündi vahel valima.

  \subsection*{Alamülesanded $K \le 4$, $K \le 8$, $K \le 16$}

  Oma ``olekut'' saame kirjeldada kolmikuga $(x, y, t)$, kus $x$ ja $y$ on ruudu
  koordinaadid, kus me asume, ning $t$ on tehtud kangitõmmete arv. Ülal veendusime,
  et pärast teatud arv kordi kangi tõmbamisi jõuab kõik algseisu tagasi. Tähistame
  seda kordade arvu $P$. $P$ väärtuseks võib võtta:
  \begin{xitem}
  \item faktoriaali $K! = 1 \cdot 2 \cdot 3 \cdots K$;
  \item vähima ühiskordse $\text{VÜK}(1, 2, \ldots, K)$;
  \item antud permutatsiooni tsüklite pikkuste vähima ühiskordse.
  \end{xitem}
  Nimekirjas allpool olevad arvud on väikesemad ja seega paremad.

  Niisiis võime olekuid $(x, y, 0)$ ja $(x, y, P)$ lugeda samaväärseteks.
  Oma tegevuse käigus saame olekust $(x, y, t)$ liikuda:
  \begin{xitem}
  \item olekutesse $(x + 1, y, t)$, $(x - 1, y, t)$, $(x, y + 1, t)$ ja $(x, y - 1, t)$,
    aga mitte siis, kui see vastab blokeeritud ruudule;
  \item kui ruudul $(x, y)$ on portaal, mis viib ruudule $(x', y')$,
    siis olekusse $(x', y', (t + 1) \bmod{P})$.
  \end{xitem}
  Leiame tavalise BFS või DFS abil, millistesse olekutesse on võimalik jõuda.

  Kui ruudul $(x, y)$ ei ole portaali, aga on münt, siis saame mündi üles korjata juhul,
  kui on võimalik jõuda vähemalt ühte olekutest $(x, y, 0)$, $(x, y, 1)$, $(x, y, 2)$ jne.
  Lisaks peame kontrollima, milliseid portaalide juures olevaid münte on võimalik üles korjata:
  portaali $i$ juures olevat münti on võimalik üles korjata siis, kui on võimalik jõuda
  ühte olekutest $(x_i, y_i, 0)$, $(x_{A_i}, y_{A_i}, 1)$, $(x_{A_{A_i}}, y_{A_{A_i}}, 2)$ jne.

  Niisugune lahendus peab üle kontrollima $N \cdot M \cdot P$ olekut. Sõltuvalt sellest, millist $P$
  väärtust kasutati, võib see teenida 30 kuni 55 punkti.

  \subsection*{Täislahendus}

  Joonistame suunatud graafi järgneval viisil: teeme tipu iga blokeerimata ruudu
  kohta. Siin dokumendis kasutame mugavuse jaoks indekseid, kus portaal $i$ asub alati
  tipus $i$. Tippu, kust alustame, tähistame $s$.
  Kui $u$ ja $v$ on naaberruudud, siis lisame servad $u \edge{0} v$
  ja $v \edge{0} u$. Lisaks lisame iga portaaliga ruudu $i$ korral servad
  $i \edge{1} A_i$ ja $A_i \edge{-1} i$.\footnote{Lahendust programmeerides võib
    serva $A_i \edge{-1} i$ lisamise ära jätta. Jääb lugejale mõtlemiseks, miks
    lahendus endiselt korrektne on.} Saadud graafil on omadus: igale servale leidub vastupidise
  suuna ja kaaluga serv. Seega ka igale teekonnale leidub vastupidise suuna ja kaaluga
  teekond.

  Nüüd saab ruudustikus liikuda ruudust $u$ ruutu $v$ täpselt $k$ kangitõmbega
  parajasti siis, kui graafis leidub teekond $u \leadsto v$ kaaluga $k$.
  Viskame graafist välja kõik tipud, kuhu algtipust $s$ teekonda ei leidu.
  Uurime, missuguste kaaludega teekonnad $s \leadsto u$ graafis leiduda saavad.

  Olgu $x$ minimaalse positiivse kaaluga tsükli kaal selles graafis (see tsükkel võib sama
  tippu või isegi serva mitu korda külastada). Nimetame sellist tsüklit $x$-tsükliks. Kui graafis ei ole
  ühtegi positiivse kaaluga tsüklit, tähendab see, et algruudust ei ole võimalik
  mööda blokeerimata ruute liikuda ühegi portaali $i$ juurde, mille korral $A_i \ne i$.
  Võimalik, et ei ole üldse võimalik ühegi portaali juurde liikuda. Sel juhul on
  ülesanne klassikaline.\footnote{Tegelikult see juht erilist eraldi käsitlemist ei vaja.
    Saab näidata, et alljärgnev lahendus lahendab ka selle juhu. Tuleb ainult
    (nulliga jagamise vältimiseks) $D[i] \equiv 0 \pmod{x}$ kontrollimiseks mitte
    teha \texttt{D[i] \% x}.}

  Näitame, et ülejäänud juhtudel:
  \begin{xitem}
  \item kui leidub teekond $s \leadsto u$ kaaluga $w$, siis leiduvad ka teekonnad
    kaaluga $w + x$ ja $w - x$;
  \item kui graafis leidub tsükkel kaaluga $w$, siis $x \mid w$ (arv $w$ jagub arvuga $x$);
  \item kui leidub teekond $s \leadsto u$ kaaluga $w$, siis ei leidu teekonda
    kaaluga $y$, kus $y \not\equiv w \pmod{x}$.
  \end{xitem}

  Esimese väite tõestuseks valime teekonna $s \leadsto u$ kaaluga $w$, liigume
  sealt $x$-tsükli juurde, teeme ühe ringi ja liigume tagasi tippu $u$.
  Nii saame teekonna $s \leadsto u$ kaaluga $w + x$. Kui liikuda mööda $x$-tsüklit vastupidises
  suunas (igale servale leidub ju vastupidise suuna ja kaaluga serv), siis saame hoopis
  teekonna kaaluga $w - x$.

  \begin{center}
    \begin{tikzpicture}
      \begin{scope}[every node/.style={circle, draw, minimum size=16pt}]
        \node (1) at (0, 0) {$s$};
        \node (2) at (8, 0) {$u$};
      \end{scope}
      
      \tikzset{decoration={snake,amplitude=.6mm,segment length=4mm,
          post length=1mm,pre length=1mm}}
      \begin{scope}[>={Stealth[black]},
          every node/.style={fill=white,circle}]
        \draw[blue, decorate, ->] (1) to[out=-5, in=-175] (2);
        \node at (4, -0.25) {{\color{blue} $w$}};

        \draw[red, decorate, ->] (4, 2) to[out=0, in=-90] (6, 4);
        \draw[red, decorate, ->] (6, 4) to[out=90, in=0] (4, 6);
        \draw[red, decorate, ->] (4, 6) to[out=180, in=90] (2, 4);
        \draw[red, decorate, ->] (2, 4) to[out=-90, in=180] (4, 2);
        \node at (4, 4) {{\color{red} $x$}};

        \draw[magenta, decorate, ->] (1) to[out=5, in=175] (2);
        \draw[magenta, decorate, ->] (2) -- (6.2, 4.2);
        \draw[magenta, decorate, ->] (6.2, 4.2) to[out=90, in=0] (4, 6.2);
        \draw[magenta, decorate, ->] (4, 6.2) to[out=180, in=90] (1.8, 4);
        \draw[magenta, decorate, ->] (1.8, 4) to[out=-90, in=180] (4, 1.8);
        \draw[magenta, decorate, ->] (4, 1.8) to[out=0, in=-90] (6.2, 3.8);
        \draw[magenta, decorate, ->] (6.2, 3.8) -- (2);
        \node at (4, 1) {{\color{magenta} $w + x$}};
      \end{scope}
    \end{tikzpicture}
  \end{center}

  Teise väite tõestuseks võtame tsükli kaaluga $w$, teeme seal ringi, liigume
  $x$-tsükli juurde, teeme seal ühe ringi ja siis liigume mööda tuldud teed tagasi.
  Saame tsükli, mille kaal on $w + x$. Analoogiliselt võime leida tsükli, mille kaal on
  $w - x$. Seega saame tsükli kaalu $x$ kaupa vähendada ja suurendada. Kui $w$ ei jaguks
  $x$-ga, siis saaksime nii tekitada tsükli, mille kaal jääks 0 ja $x$ vahele: vastuolu
  eeldusega, et $x$-tsükkel on vähima positiivse kaaluga tsükkel.

  Kolmanda väite tõestuseks võtame kaks teekonda $s \leadsto u$ kaaludega $w$ ja $y$. Kui teine
  teekond ümber pöörata ja need kaks teekonda otsapidi kokku panna, saame tsükli, mille kaal
  on $w - y$. Seega vastavalt eelmisele väitele $x \mid w - y$ ehk $w \equiv y \pmod{x}$.

  \begin{center}
    \begin{tikzpicture}
      \begin{scope}[every node/.style={circle, draw, minimum size=16pt}]
        \node (1) at (0, 0) {$s$};
        \node (2) at (8, 0) {$u$};
      \end{scope}
      
      \tikzset{decoration={snake,amplitude=.6mm,segment length=4mm,
          post length=1mm,pre length=1mm}}

      \begin{scope}[>={Stealth[black]},
          every node/.style={fill=white,circle}]
        \draw[blue, decorate, ->] (1) to[out=-25, in=-155] (2);
        \node at (4, -1.2) {{\color{blue} $y$}};

        \draw[green, decorate, ->] (1) to[out=25, in=155] (2);
        \node at (4, 1.2) {{\color{green} $w$}};

        \draw[magenta, decorate, ->] (1) to[out=15, in=165] (2);
        \draw[magenta, decorate, ->] (2) to[out=-165, in=-15] (1);
        \node at (4, 0) {{\color{magenta} $w - y$}};
      \end{scope}
    \end{tikzpicture}
  \end{center}

  Kuidas seda $x$ aga leida? Leiame esiteks (näiteks laiuti või sügavuti läbimise teel)
  iga $u$ jaoks mingi teekonna $s \leadsto u$ kaalu $D[u]$. Kui tippude $u$ ja $v$ vahel
  leidub serv kaaluga $w$, siis on graafis tsükkel kaaluga $D[u] + w - D[v]$. Järelikult
  $x \mid D[u] + w - D[v]$. Leiame kõikide selliste arvude $D[u] + w - D[v]$ suurima
  ühisteguri, tähistame selle $g$. Et $x$ on kõikide $D[u] + w - D[v]$
  jagaja, siis on ta ka nende suurima ühisteguri $g$ jagaja, seega $x \mid g$.

  Näitame, et teisest küljest $g \mid x$. Olgu minimaalse positiivse kaaluga tsükli
  tipud $u_1, u_2, \ldots, u_k$, kus $u_1 = u_k$, ja nendevaheliste servade kaalud
  $w_1, w_2, \ldots, w_k$. Teame, et
  \begin{align*}
    D[u_1] + w_1 &\equiv D[u_2] \pmod{g}, \\
    D[u_2] + w_2 &\equiv D[u_3] \pmod{g}, \\
    &\cdots \\
    D[u_{k - 1}] + w_{k - 1} &\equiv D[u_k] \pmod{g}
  \end{align*}
  Need väited kokku pannes saame, et
  \begin{align*}
    D[u_1] + w_1 + w_2 + \cdots + w_k &\equiv D[u_k] \pmod{g} \\
    D[u_1] + x &\equiv D[u_1] \pmod{g} \\
    x &\equiv 0 \pmod{g}
  \end{align*}
  ehk $g \mid x$. Järelikult $g = x$ ehk $g$ ongi selle minimaalse positiivse kaaluga
  tsükli kaal.

  Nüüd on meil iga tipu kohta teada arv $D[u]$ ja teame, et teekond $s \leadsto u$ kaaluga
  $w$ eksisteerib parajasti siis, kui $D[u] \equiv w \pmod{x}$. Kuidas selle informatsiooniga
  ülesannet lahendada? Keeruline on ainult nende müntidega, mis asuvad portaaliga ruutudel.
  Iga sellise mündi jaoks peame kontrollima, kas leidub teekond $s \leadsto i$ kaaluga 0 või
  $s \leadsto A_i$ kaaluga 1 või $s \leadsto A_{A_i}$ kaaluga 2 jne. Need väited on loomulikult
  samaväärsed: münt on võimalik üles korjata ainult juhul, kui $D[i] \equiv 0 \pmod{x}$.
\end{yl}
